\scanpage{227}
deuxième est
\[
\int_{a(s)}^{b(s)}\overline D\,d\varphi=\int_{a(s)}^{b(s)}d\,\overline D\varphi=(\overline D\varphi)(b(s))-(\overline D\varphi)(a(s)),
\]
\deleted{et les} et les deux expressions sont égales, grâce au fait que $a,b$ sont des sections horizontales. Notons aussi que ceci montre, si $a,b$ sont des sections \uline{horizontales}, qu'on aura
\[
(14)\qquad \mu(a,b)=\int_{a(s)}^{b(s)}\left(\sum\overline D_i\omega_i\right)_s=\varphi(b(s))-\varphi(a(s)),
\]
où $\varphi$ est défini par la relation (10).

Prouvons la suffisance de (7) pour la nullité de (3) en général. On va prouver ``plus \unsure{généralement|précisément}'' qu'on aura alors
\[
\sum D_{is}\int_{C_s}\omega_{is}=0
\]
lorsque $C_s$ est un 1-cycle dans $X_s$, variant différentiablement avec $s$, dont le bord est ``horizontal'' en $s$, ou ce qui revient au même
\[
\sum_i D_{is}\int_{a(s)}^{b(s)}\omega_{is}=0
\]
si $a,b$ sont des sections \uline{horizontales} de $X$ sur \deleted{\unsure{S}} \inserted{un ouvert de $S$.} \deleted{Pour} Se plaçant au voisinage d'un
\margin{226}
